Numerous research initiatives in intrusion detection, cybersecurity monitoring, and AI-driven threat intelligence have impacted the development of AI-enhanced SIEM systems. Traditional SIEM solutions, IDS/IPS frameworks, and the incorporation of machine learning for anomaly detection have all been the subject of several research efforts. This chapter examines the main current approaches, their drawbacks, and how DataLex expands on them to provide a security monitoring system that is more effective, scalable, and enhanced by AI.

\section{Related Concepts and Chapter Knowledge}
\label{sec_related_concepts}

The integration of Generative AI in cybersecurity systems, particularly in Intrusion Detection Systems (IDS), has been explored as a means to overcome the limitations of traditional SIEM and IDS platforms, which often rely on signature-based detection and predefined rules. These traditional methods typically result in high false-positive rates and limited adaptability to emerging threats, making them less effective in detecting novel or sophisticated attack patterns~\cite{juttner2023chatids}. ChatIDS addresses these issues by utilizing Generative AI to provide human-readable explanations for IDS alerts, improving the interpretability of alerts for security analysts and reducing the cognitive load associated with manual analysis. This advancement aligns with the growing need for more explainable and user-friendly security tools in modern organizations, as traditional systems often overwhelm analysts with large volumes of technical data and false positives~\cite{juttner2023chatids}.

The integration of SIEM systems with data lakes and AI/ML enhances traditional SIEM platforms, which often struggle with scalability and detecting complex threats. Traditional SIEM relies on signature-based detection, leading to high false positives and limited adaptability. By combining SIEM with data lakes, organizations can manage both structured and unstructured data more effectively, improving scalability and flexibility. The use of AI and machine learning algorithms boosts real-time detection, automates responses, and enables advanced analytics, significantly improving threat identification and response accuracy~\cite{marri2024siemdatalakes}.

The growing sophistication of cyber threats necessitates the evolution of SIEM systems beyond traditional rule-based approaches~\cite{gold2025aidriven}. AI-driven threat intelligence enhances SIEM capabilities by leveraging machine learning and real-time analytics to detect, classify, and respond to cyber threats with greater accuracy. This AI-powered integration improves SIEM's ability to identify zero-day vulnerabilities, reduce false positives, and automate incident response mechanisms. In industries like finance, healthcare, and government, AI-enhanced SIEM has shown practical benefits. However, challenges such as computational power requirements and the balance between automation and human intervention still exist. Future trends, including self-learning AI models and AI-assisted cyber forensics, are expected to further enhance real-time cybersecurity monitoring and response.

AI-driven threat intelligence is revolutionizing the field of cybersecurity by improving the capabilities of real-time threat detection, analysis, and response. Traditional security systems often rely on static, rule-based approaches that struggle to keep up with the increasingly sophisticated nature of cyber threats. The integration of machine learning models -- including supervised and unsupervised learning, reinforcement learning, and natural language processing (NLP) -- enhances the effectiveness of threat detection systems by allowing them to dynamically adapt to new attack patterns~\cite{ovabor2024aidriven}. AI frameworks are particularly valuable for addressing the growing complexity of cyber threats, enabling systems to not only detect but also predict and respond to potential security incidents. As AI tools continue to evolve, they aim to provide a proactive, dynamic approach to cybersecurity that can continuously improve its ability to combat emerging threats, overcoming many of the limitations of traditional methods.

The integration of Generative AI and Large Language Models (LLMs) into cybersecurity is enhancing real-time threat detection and response~\cite{ferrag2024genai}. Traditional security systems struggle to keep up with sophisticated threats, but LLMs like GPT-4, BERT, and LLaMA provide advanced capabilities in malware detection, intrusion prevention, and threat intelligence analysis. Despite their promise, LLMs introduce vulnerabilities such as prompt injection and adversarial instructions. Mitigation strategies, including Reinforcement Learning with Human Feedback (RLHF) and Retrieval-Augmented Generation (RAG), along with techniques like Half-Quadratic Quantization (HQQ), help optimize model performance and address challenges in training LLMs with cybersecurity datasets, aiming to improve defense mechanisms against evolving threats.

As network-based attacks continue to rise, Intrusion Detection Systems (IDS) remain crucial for detecting malicious traffic. While many Machine Learning-based IDS models exist, they often struggle with accurately detecting and classifying rare or uncommon attack types. A hybrid approach that uses synthetic data generated by Generative Adversarial Networks (GANs) to train a Deep Reinforcement Learning (DRL) model, trained on the NSL-KDD dataset for IDS evaluation, demonstrates that synthetic data generated by GANs improves the performance of the DRL model, particularly in correctly classifying minority classes, compared to training on imbalanced real-world datasets~\cite{strickland2023drlgan}.

The increasing sophistication of cybersecurity threats demands more advanced solutions. Traditional SIEM systems struggle with large data volumes and complex threats, leading to the exploration of next-generation systems that integrate data lakes and AI. Data lakes provide scalable storage for structured and unstructured data, while AI enables real-time anomaly detection, predictive analytics, and proactive threat mitigation. AI-driven SIEM systems improve threat detection speed, accuracy, and efficiency, reducing false positives and operational costs. However, challenges like data integration complexity, high costs, and privacy concerns persist. Emerging technologies such as edge AI and quantum computing are expected to further shape the future of cybersecurity~\cite{ethan2025nextgensiem}.
